% Appendix --- reproducibility details
% Compiled inline với main.tex via % Appendix --- reproducibility details
% Compiled inline với main.tex via % Appendix --- reproducibility details
% Compiled inline với main.tex via % Appendix --- reproducibility details
% Compiled inline với main.tex via \input{appendix.tex} sau \appendix

\section{Reproducibility}
\label{sec:reproducibility}

\subsection{Code and data}
Code public at \url{https://github.com/nhockid235/xling-grpo-sub3b} (Apache-2.0).
Three LoRA adapters (one per arm), all evaluation JSONs (full \texttt{responses[]}
arrays included), per-step training logs, and reward time series will be released
with the paper.

\subsection{Hyperparameters}
Full configs in \texttt{configs/}; summary below.

\paragraph{GRPO (Open-RS RS2 verbatim, except LoRA fallback).}
$\text{lr}=10^{-6}$, $G{=}6$, \texttt{max\_completion\_length}=3584,
\texttt{max\_prompt\_length}=512, $T_{\text{rollout}}=0.7$, KL $\beta{=}0.04$, 50 steps,
save every 50 steps, scheduler \texttt{cosine\_with\_min\_lr}
(\texttt{min\_lr\_rate}=0.1, \texttt{warmup\_ratio}=0.1), bf16, flash-attention 2,
gradient checkpointing. Effective batch 96 prompts (per-device 6 $\times$
\texttt{gradient\_accumulation\_steps}=16) --- Open-RS used 4$\times$A40 (per-device 6
$\times$ grad\_accum 4 $\times$ 4 GPUs); we use 1$\times$A100 with grad\_accum 16
to match.

\paragraph{LoRA configuration.} $r{=}16$, $\alpha{=}32$, dropout 0.05,
\texttt{bias=none}, \texttt{task\_type=CAUSAL\_LM}, \texttt{target\_modules}
$=\{q, k, v, o\}$\_proj. Open-RS used full-parameter; LoRA $r{=}16$ is our
single-A100 80GB fallback to fit the same effective batch size and rollout
parallelism.

\paragraph{vLLM (in-process rollout).} TRL 0.15.2 spawns vLLM 0.7.2 within the
training process. \texttt{gpu\_memory\_utilization}=0.25 (20 GB), leaving
$\sim$60 GB for training. \texttt{vllm\_max\_model\_len}=4608.

\paragraph{Reward weights.} Arms \armA{} and \armB{}: $w_{R_1}{=}w_{R_2}{=}1.0$.
Arm \armC{}: $w_{R_1}{=}w_{R_2}{=}w_{R_5}{=}1.0$ with
$R_5$'s \texttt{no\_penalty\_for\_en} flag set to \texttt{False} so $R_5$ fires on
every prompt (not just non-EN).

\paragraph{Eval.} vLLM 0.7.2, dtype bfloat16, $T{=}0$ for pass@1 (greedy), $T{=}0.7$
top\_p=0.95 for AIME-2024 maj@8, max\_tokens=8192, no early stop tokens.
Open-RS \texttt{lighteval} \texttt{MATH\_QUERY\_TEMPLATE} verbatim; no system
prompt (template embedded in user message).

\subsection{Datasets (HuggingFace IDs)}
\begin{itemize}\itemsep0pt
\item \armA{}, \armC{} training: \texttt{knoveleng/open-rs} (MIT, 7K, \texttt{split=train}).
\item \armB{} training:
\texttt{5CD-AI/Vietnamese-meta-math-MetaMathQA-40K-gg-translated}; license unset
on Hub (parent MetaMathQA = MIT). We extracted 5{,}203 records from a 7K random
subset by retaining only those whose Vietnamese response contains the
``\DJ\'ap \'an l\`a'' (Vietnamese for ``the answer is'') answer marker.
\item Eval: \texttt{HuggingFaceH4/MATH-500} (500 test), \texttt{Maxwell-Jia/AIME\_2024}
(MIT, 30 records; field names capitalized: \texttt{Problem}, \texttt{Solution},
\texttt{Answer}, \texttt{ID}), \texttt{knoveleng/AMC-23} (40 records; license
unspecified on Hub, originally AoPS wiki content).
\end{itemize}

\subsection{Decontamination}
We did not perform train/test decontamination for the present arms because the
training datasets (\texttt{knoveleng/open-rs} and \texttt{5CD-AI/Vietnamese-MetaMathQA})
are upstream of the evaluation datasets in different domains; in addition,
Open-RS RS2 itself does not decontaminate. A future ablation will quantify any
contamination effect via 8-gram overlap statistics.

\subsection{Compute and cost}
Training+eval: 1$\times$ NVIDIA A100-SXM4-80GB on Vast.ai cloud, \$1.50/h.
Total wallclock for the three arms (50 steps each + ckpt-50 evaluation on
3 benchmarks): 11 h 39 min, \$17.50.
Setup, dependency installation, and failed retries: \$2.25.
Total project spend: \$19.75 of \$180 budget.

\subsection{Random seeds}
All training: \texttt{seed=42}. AIME-2024 maj@8 uses seeds 42--49 (8 stochastic
samples per problem). The single-seed limitation is the most significant
threat to validity; we plan a 3-seed extension (seeds 42, 123, 7) for camera-ready.

\subsection{Software versions}
Python 3.12.3, \texttt{torch==2.5.1+cu124}, \texttt{transformers==4.49.0},
\texttt{trl==0.15.2}, \texttt{tokenizers==0.21.0}, \texttt{vllm==0.7.2},
\texttt{accelerate==1.2.1}, \texttt{peft==0.14.0}, \texttt{datasets==4.8.5},
\texttt{flash\_attn==2.7.2.post1}, \texttt{sympy==1.13.1},
\texttt{math\_verify==0.9.0}, \texttt{fasttext-wheel==0.9.2}
(with \texttt{lid.176.bin} 126 MB).

\textbf{Note on TRL version.} \texttt{GRPOTrainer} first appeared in TRL 0.14.0;
\texttt{GRPOConfig.reward\_weights} field requires TRL $\geq$ 0.15.0. TRL 0.15.2
spawns vLLM in-process; TRL 0.16.0+ requires an external vLLM server. Our pipeline
targets the 0.15.2 in-process model.

\textbf{Note on sympy version.} The CLAUDE.md project specification originally
pinned \texttt{sympy<1.13} citing Math-Verify breakage; in practice
\texttt{math\_verify==0.9.0} works correctly with \texttt{sympy==1.13.1}
(verified by running \texttt{verify(parse(``1/2''), parse(``0.5''))}).

\subsection{Hardware notes}
The single-A100 LoRA constraint required several config adjustments not present
in Open-RS:
\begin{itemize}\itemsep0pt
\item \texttt{vllm\_gpu\_memory\_utilization} reduced from Open-RS's 0.7 to 0.25 to
leave room for the LoRA training pass.
\item Effective batch maintained at 96 (Open-RS) by setting
\texttt{gradient\_accumulation\_steps}=16 instead of 4 (4$\times$A40) on a single
device.
\item \texttt{gradient\_checkpointing} enabled with
\texttt{use\_reentrant=False} to fit activations in remaining VRAM.
\end{itemize}

\subsection{Limitations not in main text}
\begin{itemize}\itemsep0pt
\item \texttt{maj@4} on AMC23 returns 0 across all conditions including base --- likely
a sampling-seed handling bug in our majority-vote implementation; pass@1 numbers
verified independent of this bug.
\item Eval pipeline reports a $-12.9$ pp gap to Open-RS reported AMC23 base
(50.0\% vs 62.9\%) and a $-23.4$ pp gap on MATH-500 (59.4\% vs 82.8\%) even with
verbatim Open-RS \texttt{lighteval} prompt. Suspected causes: \texttt{Math-Verify}
configuration differences between our adapter and Open-RS's
\texttt{LatexExtractionConfig(boxed\_match\_priority=0)}, or vLLM tokenizer edge
cases. We report numbers vs our own base rather than vs Open-RS reported.
\item AIME-2024 baseline is closer to Open-RS reported ($26.7$\% vs $28.8$\%, gap
$-2.1$ pp), suggesting eval pipeline issues are problem-difficulty-dependent.
\end{itemize}
 sau \appendix

\section{Reproducibility}
\label{sec:reproducibility}

\subsection{Code and data}
Code public at \url{https://github.com/nhockid235/xling-grpo-sub3b} (Apache-2.0).
Three LoRA adapters (one per arm), all evaluation JSONs (full \texttt{responses[]}
arrays included), per-step training logs, and reward time series will be released
with the paper.

\subsection{Hyperparameters}
Full configs in \texttt{configs/}; summary below.

\paragraph{GRPO (Open-RS RS2 verbatim, except LoRA fallback).}
$\text{lr}=10^{-6}$, $G{=}6$, \texttt{max\_completion\_length}=3584,
\texttt{max\_prompt\_length}=512, $T_{\text{rollout}}=0.7$, KL $\beta{=}0.04$, 50 steps,
save every 50 steps, scheduler \texttt{cosine\_with\_min\_lr}
(\texttt{min\_lr\_rate}=0.1, \texttt{warmup\_ratio}=0.1), bf16, flash-attention 2,
gradient checkpointing. Effective batch 96 prompts (per-device 6 $\times$
\texttt{gradient\_accumulation\_steps}=16) --- Open-RS used 4$\times$A40 (per-device 6
$\times$ grad\_accum 4 $\times$ 4 GPUs); we use 1$\times$A100 with grad\_accum 16
to match.

\paragraph{LoRA configuration.} $r{=}16$, $\alpha{=}32$, dropout 0.05,
\texttt{bias=none}, \texttt{task\_type=CAUSAL\_LM}, \texttt{target\_modules}
$=\{q, k, v, o\}$\_proj. Open-RS used full-parameter; LoRA $r{=}16$ is our
single-A100 80GB fallback to fit the same effective batch size and rollout
parallelism.

\paragraph{vLLM (in-process rollout).} TRL 0.15.2 spawns vLLM 0.7.2 within the
training process. \texttt{gpu\_memory\_utilization}=0.25 (20 GB), leaving
$\sim$60 GB for training. \texttt{vllm\_max\_model\_len}=4608.

\paragraph{Reward weights.} Arms \armA{} and \armB{}: $w_{R_1}{=}w_{R_2}{=}1.0$.
Arm \armC{}: $w_{R_1}{=}w_{R_2}{=}w_{R_5}{=}1.0$ with
$R_5$'s \texttt{no\_penalty\_for\_en} flag set to \texttt{False} so $R_5$ fires on
every prompt (not just non-EN).

\paragraph{Eval.} vLLM 0.7.2, dtype bfloat16, $T{=}0$ for pass@1 (greedy), $T{=}0.7$
top\_p=0.95 for AIME-2024 maj@8, max\_tokens=8192, no early stop tokens.
Open-RS \texttt{lighteval} \texttt{MATH\_QUERY\_TEMPLATE} verbatim; no system
prompt (template embedded in user message).

\subsection{Datasets (HuggingFace IDs)}
\begin{itemize}\itemsep0pt
\item \armA{}, \armC{} training: \texttt{knoveleng/open-rs} (MIT, 7K, \texttt{split=train}).
\item \armB{} training:
\texttt{5CD-AI/Vietnamese-meta-math-MetaMathQA-40K-gg-translated}; license unset
on Hub (parent MetaMathQA = MIT). We extracted 5{,}203 records from a 7K random
subset by retaining only those whose Vietnamese response contains the
``\DJ\'ap \'an l\`a'' (Vietnamese for ``the answer is'') answer marker.
\item Eval: \texttt{HuggingFaceH4/MATH-500} (500 test), \texttt{Maxwell-Jia/AIME\_2024}
(MIT, 30 records; field names capitalized: \texttt{Problem}, \texttt{Solution},
\texttt{Answer}, \texttt{ID}), \texttt{knoveleng/AMC-23} (40 records; license
unspecified on Hub, originally AoPS wiki content).
\end{itemize}

\subsection{Decontamination}
We did not perform train/test decontamination for the present arms because the
training datasets (\texttt{knoveleng/open-rs} and \texttt{5CD-AI/Vietnamese-MetaMathQA})
are upstream of the evaluation datasets in different domains; in addition,
Open-RS RS2 itself does not decontaminate. A future ablation will quantify any
contamination effect via 8-gram overlap statistics.

\subsection{Compute and cost}
Training+eval: 1$\times$ NVIDIA A100-SXM4-80GB on Vast.ai cloud, \$1.50/h.
Total wallclock for the three arms (50 steps each + ckpt-50 evaluation on
3 benchmarks): 11 h 39 min, \$17.50.
Setup, dependency installation, and failed retries: \$2.25.
Total project spend: \$19.75 of \$180 budget.

\subsection{Random seeds}
All training: \texttt{seed=42}. AIME-2024 maj@8 uses seeds 42--49 (8 stochastic
samples per problem). The single-seed limitation is the most significant
threat to validity; we plan a 3-seed extension (seeds 42, 123, 7) for camera-ready.

\subsection{Software versions}
Python 3.12.3, \texttt{torch==2.5.1+cu124}, \texttt{transformers==4.49.0},
\texttt{trl==0.15.2}, \texttt{tokenizers==0.21.0}, \texttt{vllm==0.7.2},
\texttt{accelerate==1.2.1}, \texttt{peft==0.14.0}, \texttt{datasets==4.8.5},
\texttt{flash\_attn==2.7.2.post1}, \texttt{sympy==1.13.1},
\texttt{math\_verify==0.9.0}, \texttt{fasttext-wheel==0.9.2}
(with \texttt{lid.176.bin} 126 MB).

\textbf{Note on TRL version.} \texttt{GRPOTrainer} first appeared in TRL 0.14.0;
\texttt{GRPOConfig.reward\_weights} field requires TRL $\geq$ 0.15.0. TRL 0.15.2
spawns vLLM in-process; TRL 0.16.0+ requires an external vLLM server. Our pipeline
targets the 0.15.2 in-process model.

\textbf{Note on sympy version.} The CLAUDE.md project specification originally
pinned \texttt{sympy<1.13} citing Math-Verify breakage; in practice
\texttt{math\_verify==0.9.0} works correctly with \texttt{sympy==1.13.1}
(verified by running \texttt{verify(parse(``1/2''), parse(``0.5''))}).

\subsection{Hardware notes}
The single-A100 LoRA constraint required several config adjustments not present
in Open-RS:
\begin{itemize}\itemsep0pt
\item \texttt{vllm\_gpu\_memory\_utilization} reduced from Open-RS's 0.7 to 0.25 to
leave room for the LoRA training pass.
\item Effective batch maintained at 96 (Open-RS) by setting
\texttt{gradient\_accumulation\_steps}=16 instead of 4 (4$\times$A40) on a single
device.
\item \texttt{gradient\_checkpointing} enabled with
\texttt{use\_reentrant=False} to fit activations in remaining VRAM.
\end{itemize}

\subsection{Limitations not in main text}
\begin{itemize}\itemsep0pt
\item \texttt{maj@4} on AMC23 returns 0 across all conditions including base --- likely
a sampling-seed handling bug in our majority-vote implementation; pass@1 numbers
verified independent of this bug.
\item Eval pipeline reports a $-12.9$ pp gap to Open-RS reported AMC23 base
(50.0\% vs 62.9\%) and a $-23.4$ pp gap on MATH-500 (59.4\% vs 82.8\%) even with
verbatim Open-RS \texttt{lighteval} prompt. Suspected causes: \texttt{Math-Verify}
configuration differences between our adapter and Open-RS's
\texttt{LatexExtractionConfig(boxed\_match\_priority=0)}, or vLLM tokenizer edge
cases. We report numbers vs our own base rather than vs Open-RS reported.
\item AIME-2024 baseline is closer to Open-RS reported ($26.7$\% vs $28.8$\%, gap
$-2.1$ pp), suggesting eval pipeline issues are problem-difficulty-dependent.
\end{itemize}
 sau \appendix

\section{Reproducibility}
\label{sec:reproducibility}

\subsection{Code and data}
Code public at \url{https://github.com/nhockid235/xling-grpo-sub3b} (Apache-2.0).
Three LoRA adapters (one per arm), all evaluation JSONs (full \texttt{responses[]}
arrays included), per-step training logs, and reward time series will be released
with the paper.

\subsection{Hyperparameters}
Full configs in \texttt{configs/}; summary below.

\paragraph{GRPO (Open-RS RS2 verbatim, except LoRA fallback).}
$\text{lr}=10^{-6}$, $G{=}6$, \texttt{max\_completion\_length}=3584,
\texttt{max\_prompt\_length}=512, $T_{\text{rollout}}=0.7$, KL $\beta{=}0.04$, 50 steps,
save every 50 steps, scheduler \texttt{cosine\_with\_min\_lr}
(\texttt{min\_lr\_rate}=0.1, \texttt{warmup\_ratio}=0.1), bf16, flash-attention 2,
gradient checkpointing. Effective batch 96 prompts (per-device 6 $\times$
\texttt{gradient\_accumulation\_steps}=16) --- Open-RS used 4$\times$A40 (per-device 6
$\times$ grad\_accum 4 $\times$ 4 GPUs); we use 1$\times$A100 with grad\_accum 16
to match.

\paragraph{LoRA configuration.} $r{=}16$, $\alpha{=}32$, dropout 0.05,
\texttt{bias=none}, \texttt{task\_type=CAUSAL\_LM}, \texttt{target\_modules}
$=\{q, k, v, o\}$\_proj. Open-RS used full-parameter; LoRA $r{=}16$ is our
single-A100 80GB fallback to fit the same effective batch size and rollout
parallelism.

\paragraph{vLLM (in-process rollout).} TRL 0.15.2 spawns vLLM 0.7.2 within the
training process. \texttt{gpu\_memory\_utilization}=0.25 (20 GB), leaving
$\sim$60 GB for training. \texttt{vllm\_max\_model\_len}=4608.

\paragraph{Reward weights.} Arms \armA{} and \armB{}: $w_{R_1}{=}w_{R_2}{=}1.0$.
Arm \armC{}: $w_{R_1}{=}w_{R_2}{=}w_{R_5}{=}1.0$ with
$R_5$'s \texttt{no\_penalty\_for\_en} flag set to \texttt{False} so $R_5$ fires on
every prompt (not just non-EN).

\paragraph{Eval.} vLLM 0.7.2, dtype bfloat16, $T{=}0$ for pass@1 (greedy), $T{=}0.7$
top\_p=0.95 for AIME-2024 maj@8, max\_tokens=8192, no early stop tokens.
Open-RS \texttt{lighteval} \texttt{MATH\_QUERY\_TEMPLATE} verbatim; no system
prompt (template embedded in user message).

\subsection{Datasets (HuggingFace IDs)}
\begin{itemize}\itemsep0pt
\item \armA{}, \armC{} training: \texttt{knoveleng/open-rs} (MIT, 7K, \texttt{split=train}).
\item \armB{} training:
\texttt{5CD-AI/Vietnamese-meta-math-MetaMathQA-40K-gg-translated}; license unset
on Hub (parent MetaMathQA = MIT). We extracted 5{,}203 records from a 7K random
subset by retaining only those whose Vietnamese response contains the
``\DJ\'ap \'an l\`a'' (Vietnamese for ``the answer is'') answer marker.
\item Eval: \texttt{HuggingFaceH4/MATH-500} (500 test), \texttt{Maxwell-Jia/AIME\_2024}
(MIT, 30 records; field names capitalized: \texttt{Problem}, \texttt{Solution},
\texttt{Answer}, \texttt{ID}), \texttt{knoveleng/AMC-23} (40 records; license
unspecified on Hub, originally AoPS wiki content).
\end{itemize}

\subsection{Decontamination}
We did not perform train/test decontamination for the present arms because the
training datasets (\texttt{knoveleng/open-rs} and \texttt{5CD-AI/Vietnamese-MetaMathQA})
are upstream of the evaluation datasets in different domains; in addition,
Open-RS RS2 itself does not decontaminate. A future ablation will quantify any
contamination effect via 8-gram overlap statistics.

\subsection{Compute and cost}
Training+eval: 1$\times$ NVIDIA A100-SXM4-80GB on Vast.ai cloud, \$1.50/h.
Total wallclock for the three arms (50 steps each + ckpt-50 evaluation on
3 benchmarks): 11 h 39 min, \$17.50.
Setup, dependency installation, and failed retries: \$2.25.
Total project spend: \$19.75 of \$180 budget.

\subsection{Random seeds}
All training: \texttt{seed=42}. AIME-2024 maj@8 uses seeds 42--49 (8 stochastic
samples per problem). The single-seed limitation is the most significant
threat to validity; we plan a 3-seed extension (seeds 42, 123, 7) for camera-ready.

\subsection{Software versions}
Python 3.12.3, \texttt{torch==2.5.1+cu124}, \texttt{transformers==4.49.0},
\texttt{trl==0.15.2}, \texttt{tokenizers==0.21.0}, \texttt{vllm==0.7.2},
\texttt{accelerate==1.2.1}, \texttt{peft==0.14.0}, \texttt{datasets==4.8.5},
\texttt{flash\_attn==2.7.2.post1}, \texttt{sympy==1.13.1},
\texttt{math\_verify==0.9.0}, \texttt{fasttext-wheel==0.9.2}
(with \texttt{lid.176.bin} 126 MB).

\textbf{Note on TRL version.} \texttt{GRPOTrainer} first appeared in TRL 0.14.0;
\texttt{GRPOConfig.reward\_weights} field requires TRL $\geq$ 0.15.0. TRL 0.15.2
spawns vLLM in-process; TRL 0.16.0+ requires an external vLLM server. Our pipeline
targets the 0.15.2 in-process model.

\textbf{Note on sympy version.} The CLAUDE.md project specification originally
pinned \texttt{sympy<1.13} citing Math-Verify breakage; in practice
\texttt{math\_verify==0.9.0} works correctly with \texttt{sympy==1.13.1}
(verified by running \texttt{verify(parse(``1/2''), parse(``0.5''))}).

\subsection{Hardware notes}
The single-A100 LoRA constraint required several config adjustments not present
in Open-RS:
\begin{itemize}\itemsep0pt
\item \texttt{vllm\_gpu\_memory\_utilization} reduced from Open-RS's 0.7 to 0.25 to
leave room for the LoRA training pass.
\item Effective batch maintained at 96 (Open-RS) by setting
\texttt{gradient\_accumulation\_steps}=16 instead of 4 (4$\times$A40) on a single
device.
\item \texttt{gradient\_checkpointing} enabled with
\texttt{use\_reentrant=False} to fit activations in remaining VRAM.
\end{itemize}

\subsection{Limitations not in main text}
\begin{itemize}\itemsep0pt
\item \texttt{maj@4} on AMC23 returns 0 across all conditions including base --- likely
a sampling-seed handling bug in our majority-vote implementation; pass@1 numbers
verified independent of this bug.
\item Eval pipeline reports a $-12.9$ pp gap to Open-RS reported AMC23 base
(50.0\% vs 62.9\%) and a $-23.4$ pp gap on MATH-500 (59.4\% vs 82.8\%) even with
verbatim Open-RS \texttt{lighteval} prompt. Suspected causes: \texttt{Math-Verify}
configuration differences between our adapter and Open-RS's
\texttt{LatexExtractionConfig(boxed\_match\_priority=0)}, or vLLM tokenizer edge
cases. We report numbers vs our own base rather than vs Open-RS reported.
\item AIME-2024 baseline is closer to Open-RS reported ($26.7$\% vs $28.8$\%, gap
$-2.1$ pp), suggesting eval pipeline issues are problem-difficulty-dependent.
\end{itemize}
 sau \appendix

\section{Reproducibility}
\label{sec:reproducibility}

\subsection{Code and data}
Code public at \url{https://github.com/nhockid235/xling-grpo-sub3b} (Apache-2.0).
Three LoRA adapters (one per arm), all evaluation JSONs (full \texttt{responses[]}
arrays included), per-step training logs, and reward time series will be released
with the paper.

\subsection{Hyperparameters}
Full configs in \texttt{configs/}; summary below.

\paragraph{GRPO (Open-RS RS2 verbatim, except LoRA fallback).}
$\text{lr}=10^{-6}$, $G{=}6$, \texttt{max\_completion\_length}=3584,
\texttt{max\_prompt\_length}=512, $T_{\text{rollout}}=0.7$, KL $\beta{=}0.04$, 50 steps,
save every 50 steps, scheduler \texttt{cosine\_with\_min\_lr}
(\texttt{min\_lr\_rate}=0.1, \texttt{warmup\_ratio}=0.1), bf16, flash-attention 2,
gradient checkpointing. Effective batch 96 prompts (per-device 6 $\times$
\texttt{gradient\_accumulation\_steps}=16) --- Open-RS used 4$\times$A40 (per-device 6
$\times$ grad\_accum 4 $\times$ 4 GPUs); we use 1$\times$A100 with grad\_accum 16
to match.

\paragraph{LoRA configuration.} $r{=}16$, $\alpha{=}32$, dropout 0.05,
\texttt{bias=none}, \texttt{task\_type=CAUSAL\_LM}, \texttt{target\_modules}
$=\{q, k, v, o\}$\_proj. Open-RS used full-parameter; LoRA $r{=}16$ is our
single-A100 80GB fallback to fit the same effective batch size and rollout
parallelism.

\paragraph{vLLM (in-process rollout).} TRL 0.15.2 spawns vLLM 0.7.2 within the
training process. \texttt{gpu\_memory\_utilization}=0.25 (20 GB), leaving
$\sim$60 GB for training. \texttt{vllm\_max\_model\_len}=4608.

\paragraph{Reward weights.} Arms \armA{} and \armB{}: $w_{R_1}{=}w_{R_2}{=}1.0$.
Arm \armC{}: $w_{R_1}{=}w_{R_2}{=}w_{R_5}{=}1.0$ with
$R_5$'s \texttt{no\_penalty\_for\_en} flag set to \texttt{False} so $R_5$ fires on
every prompt (not just non-EN).

\paragraph{Eval.} vLLM 0.7.2, dtype bfloat16, $T{=}0$ for pass@1 (greedy), $T{=}0.7$
top\_p=0.95 for AIME-2024 maj@8, max\_tokens=8192, no early stop tokens.
Open-RS \texttt{lighteval} \texttt{MATH\_QUERY\_TEMPLATE} verbatim; no system
prompt (template embedded in user message).

\subsection{Datasets (HuggingFace IDs)}
\begin{itemize}\itemsep0pt
\item \armA{}, \armC{} training: \texttt{knoveleng/open-rs} (MIT, 7K, \texttt{split=train}).
\item \armB{} training:
\texttt{5CD-AI/Vietnamese-meta-math-MetaMathQA-40K-gg-translated}; license unset
on Hub (parent MetaMathQA = MIT). We extracted 5{,}203 records from a 7K random
subset by retaining only those whose Vietnamese response contains the
``\DJ\'ap \'an l\`a'' (Vietnamese for ``the answer is'') answer marker.
\item Eval: \texttt{HuggingFaceH4/MATH-500} (500 test), \texttt{Maxwell-Jia/AIME\_2024}
(MIT, 30 records; field names capitalized: \texttt{Problem}, \texttt{Solution},
\texttt{Answer}, \texttt{ID}), \texttt{knoveleng/AMC-23} (40 records; license
unspecified on Hub, originally AoPS wiki content).
\end{itemize}

\subsection{Decontamination}
We did not perform train/test decontamination for the present arms because the
training datasets (\texttt{knoveleng/open-rs} and \texttt{5CD-AI/Vietnamese-MetaMathQA})
are upstream of the evaluation datasets in different domains; in addition,
Open-RS RS2 itself does not decontaminate. A future ablation will quantify any
contamination effect via 8-gram overlap statistics.

\subsection{Compute and cost}
Training+eval: 1$\times$ NVIDIA A100-SXM4-80GB on Vast.ai cloud, \$1.50/h.
Total wallclock for the three arms (50 steps each + ckpt-50 evaluation on
3 benchmarks): 11 h 39 min, \$17.50.
Setup, dependency installation, and failed retries: \$2.25.
Total project spend: \$19.75 of \$180 budget.

\subsection{Random seeds}
All training: \texttt{seed=42}. AIME-2024 maj@8 uses seeds 42--49 (8 stochastic
samples per problem). The single-seed limitation is the most significant
threat to validity; we plan a 3-seed extension (seeds 42, 123, 7) for camera-ready.

\subsection{Software versions}
Python 3.12.3, \texttt{torch==2.5.1+cu124}, \texttt{transformers==4.49.0},
\texttt{trl==0.15.2}, \texttt{tokenizers==0.21.0}, \texttt{vllm==0.7.2},
\texttt{accelerate==1.2.1}, \texttt{peft==0.14.0}, \texttt{datasets==4.8.5},
\texttt{flash\_attn==2.7.2.post1}, \texttt{sympy==1.13.1},
\texttt{math\_verify==0.9.0}, \texttt{fasttext-wheel==0.9.2}
(with \texttt{lid.176.bin} 126 MB).

\textbf{Note on TRL version.} \texttt{GRPOTrainer} first appeared in TRL 0.14.0;
\texttt{GRPOConfig.reward\_weights} field requires TRL $\geq$ 0.15.0. TRL 0.15.2
spawns vLLM in-process; TRL 0.16.0+ requires an external vLLM server. Our pipeline
targets the 0.15.2 in-process model.

\textbf{Note on sympy version.} The CLAUDE.md project specification originally
pinned \texttt{sympy<1.13} citing Math-Verify breakage; in practice
\texttt{math\_verify==0.9.0} works correctly with \texttt{sympy==1.13.1}
(verified by running \texttt{verify(parse(``1/2''), parse(``0.5''))}).

\subsection{Hardware notes}
The single-A100 LoRA constraint required several config adjustments not present
in Open-RS:
\begin{itemize}\itemsep0pt
\item \texttt{vllm\_gpu\_memory\_utilization} reduced from Open-RS's 0.7 to 0.25 to
leave room for the LoRA training pass.
\item Effective batch maintained at 96 (Open-RS) by setting
\texttt{gradient\_accumulation\_steps}=16 instead of 4 (4$\times$A40) on a single
device.
\item \texttt{gradient\_checkpointing} enabled with
\texttt{use\_reentrant=False} to fit activations in remaining VRAM.
\end{itemize}

\subsection{Limitations not in main text}
\begin{itemize}\itemsep0pt
\item \texttt{maj@4} on AMC23 returns 0 across all conditions including base --- likely
a sampling-seed handling bug in our majority-vote implementation; pass@1 numbers
verified independent of this bug.
\item Eval pipeline reports a $-12.9$ pp gap to Open-RS reported AMC23 base
(50.0\% vs 62.9\%) and a $-23.4$ pp gap on MATH-500 (59.4\% vs 82.8\%) even with
verbatim Open-RS \texttt{lighteval} prompt. Suspected causes: \texttt{Math-Verify}
configuration differences between our adapter and Open-RS's
\texttt{LatexExtractionConfig(boxed\_match\_priority=0)}, or vLLM tokenizer edge
cases. We report numbers vs our own base rather than vs Open-RS reported.
\item AIME-2024 baseline is closer to Open-RS reported ($26.7$\% vs $28.8$\%, gap
$-2.1$ pp), suggesting eval pipeline issues are problem-difficulty-dependent.
\end{itemize}
